%==============================================================================
%  06-why-ai-fails.tex — Section 5: Why AI Fails
%==============================================================================
\strategyPart{18 \textbar\ Why AI Fails}{The Organizational Gap}

The most important insight in the AI research literature is that failures are
predictable, patterned, and organizational — not random technical accidents.
Naming the failure modes is the first step to designing them out.

\section{The evidence on failure}

\begin{itemize}
  \item Gartner warned that through 2022, \textbf{85\% of AI projects would
        deliver erroneous outcomes} due to bias in data, algorithms, or the
        responsible AI teams managing them [Gartner 2018].
  \item Gartner further predicts that through 2025, \textbf{30\% of generative
        AI projects will be abandoned after proof of concept} — the point where
        value must move from demo to production [Gartner 2023].
  \item Seven in ten companies report \textbf{minimal or no impact} from AI so
        far, and 40\% of significant AI investors report no bottom-line impact
        [MIT SMR--BCG 2019].
\end{itemize}

The pattern is consistent: the bottleneck sits between proof of concept and
production, and it is organizational.

\section{The four failure modes}

\vspace{-0.3em}
\infobox{\textbf{1 · The technology-first fallacy}}{%
  Capability is acquired before a problem is defined. The result is a portfolio
  of solutions in search of problems — demos that never become P\&L impact
  [MIT SMR--BCG 2019].}

\infobox{\textbf{2 · The data trust deficit}}{%
  Models are only as trustworthy as the data beneath them. Unmanaged data
  quality, silos and weak lineage produce models that no business owner will
  act on [Gartner 2021].}

\infobox{\textbf{3 · Cultural resistance \& low literacy}}{%
  When employees fear substitution and cannot interrogate AI outputs, they
  bypass the system. Augmentation without literacy produces adoption failure
  [HBR 2018].}

\infobox{\textbf{4 · Poor problem definition}}{%
  Fuzzy problems produce fuzzy metrics. Without a defined problem, a KPI, and
  an owner, an AI project cannot fail visibly — and therefore cannot be fixed
  or scaled [McKinsey 2023].}

\section{Why the gap is widening}

The agentic frontier is making the organizational gap more expensive, not less.
As Gartner projects that one-third of enterprise software will include agentic
AI by 2028 and 15\% of day-to-day decisions will become autonomous [Gartner
2024], the cost of weak governance, untrusted data and unprepared workforces
scales with the autonomy of the system. Organizations that close the
organizational gap first will compound the advantage; those that do not will
automate their errors faster.

\takeaway{AI failures are predictable and organizational; prevent them with problem definition, data trust, literacy and governance.}
